\section{Protocole méthodologique et architecture du projet}
\label{sec:methodology}
\subsection{Trois objets scientifiques distincts}
La prévision du prix ajusté de NVDA, les interdépendances conditionnelles des
rendements de NVDA, AMD et SOXX et la variance des innovations constituent trois objets
différents. La qualité d'un diagnostic de moyenne ou de variance n'est pas un
gain de prévision de prix. Le schéma de la figure~\ref{fig:pipeline} situe ces
analyses et leurs livrables ; GARCH reste hors du classement final des prix.

\subsection{Temporalité, validation interne et holdout}
Le dataset est partagé chronologiquement avant la sélection des modèles :
1980 lignes TRAIN puis 220 lignes TEST. Les décisions de transformation et les
diagnostics de modélisation reposent sur TRAIN. Les figures descriptives couvrant
toute la période sont identifiées comme telles ; leur existence n'est pas
présentée comme un accès anticipé du nouvel évaluateur final.

ARIMA utilise une validation interne rolling-origin ; les réseaux utilisent
trois folds expanding-window et une validation d'arrêt interne distincte des
blocs externes. Les scalers sont ajustés dans les seuls préfixes permis.
L'observation connue à une nouvelle origine peut enrichir la fenêtre passée,
sans autoriser un nouveau fit sur le holdout.

Avant le scoring final, les spécifications, paramètres et règles sont gelés.
Tous les modèles de prix sont comparés sur les mêmes 211 origines à cinq et dix
séances observées. Les coefficients/poids/scalers restent fixes ; seule
l'information déjà observée peut avancer. Le benchmark est
\begin{equation}
\label{eq:naive}
\widehat P_{t+h\mid t}^{\,\mathrm{Naive}}=P_t,\qquad h\in\{5,10\}.
\end{equation}
Son rôle de référence ne prouve ni bruit blanc des innovations ni optimalité
universelle. Le lock Naive interne de Phase 3B, le challenger ARIMA(2,1,2)
pré-déclaré et la recommandation finale sont trois décisions distinctes.

% Documentary process diagram; no model execution and no invented result.
\begin{figure}[!htbp]
\centering
\begin{tikzpicture}[
  font=\small,
  box/.style={draw=academicblue,rounded corners=2pt,fill=lightblue,
    align=center,text width=14.8cm,inner sep=6pt},
  model/.style={box,text width=4.2cm,minimum height=2.4cm},
  flow/.style={-{Latex[length=2.5mm]},draw=academicblue,thick},
  risk/.style={flow,draw=riskgold,dashed}]
\node[box] (input) {\textbf{Entrées et provenance}\\
  NVDA, AMD, SOXX ; Yahoo/yfinance ; snapshots originaux et SHA256};
\node[box,below=7mm of input] (data) {\textbf{Préparation et alignement}\\
  Calendriers communs ; log-prix, rendements et volatilité passée\\
  Séparation chronologique : TRAIN 90\,\% / TEST 10\,\%};
\node[box,below=7mm of data] (eda) {\textbf{EDA et stationnarité}\\
  Description distinguée des décisions TRAIN ; ACF/PACF, ADF/KPSS\\
  Protocoles scientifiques gelés avant évaluation finale};
\node[model,anchor=north] (classic) at ($(eda.south)+(-5,-0.9)$)
  {\textbf{Moyenne et interactions}\\
  ARIMA et validation\\VAR(9), Granger, Johansen\\
  Rang 0 : VECM non justifié};
\node[model,anchor=north] (neural) at ($(eda.south)+(0,-0.9)$)
  {\textbf{Réseaux récurrents}\\
  LSTM C / GRU C\\Folds et scalers TRAIN-only\\18 fits + 2 refits};
\node[model,anchor=north,draw=riskgold,fill=white] (garch)
  at ($(eda.south)+(5,-0.9)$)
  {\textbf{Variance conditionnelle}\\
  ARCH-LM puis GARCH(1,1)\\Normal verrouillé ; sigma\\
  Hors classement des prix};
\node[box,text width=9.6cm,anchor=north] (test)
  at ($(classic.south)!0.5!(neural.south)+(0,-0.7)$)
  {\textbf{Évaluation finale des prix sur holdout}\\
  Gate ; modèles gelés ; benchmark Naive\\
  211 origines communes ; h5/h10\\MSE/RMSE/MAE/MAPE ; DM/Holm};
\node[box,text width=9.6cm,below=7mm of test] (naive)
  {\textbf{Recommandation finale : Naive}\\
  Aucun challenger ne satisfait l'ensemble de la règle pré-déclarée};
\node[box,anchor=north] (llm) at ($(naive.south)+(2.5,-0.7)$)
  {\textbf{Analyse LLM rétrospective}\\
  Réponse officielle conservée ; comparaison critique\\
  Audit documentaire de 35 affirmations};
\node[box,below=7mm of llm] (delivery)
  {\textbf{Reproduction et livrables}\\
  FAST : artefacts gelés ; FULL : run isolé exécuté, §8.3 PARTIAL\\
  Notebook FULL et machine vierge à démontrer\\
  Rapports Word/PDF conservés et projet LaTeX séparé};
\draw[flow] (input.south) -- (data.north);
\draw[flow] (data.south) -- (eda.north);
\draw[flow] (eda.south) -- ++(0,-0.35) -| (classic.north);
\draw[flow] (eda.south) -- (neural.north);
\draw[risk] (eda.south) -- ++(0,-0.35) -| (garch.north);
\draw[flow] (classic.south) -- ([xshift=-2.5cm]test.north);
\draw[flow] (neural.south) -- ([xshift=2.5cm]test.north);
\draw[flow] (test.south) -- (naive.north);
\draw[flow] (naive.south) -- ([xshift=-2.5cm]llm.north);
\draw[risk] (garch.south) -- (garch.south |- llm.north);
\draw[flow] (llm.south) -- (delivery.north);
\end{tikzpicture}
\caption{Architecture documentaire et processus scientifique du projet.}
\label{fig:pipeline}
\sourcecaption{Synthèse des protocoles et audits Phases 1--9C.
  Traits pleins : chaîne des prévisions de prix ; pointillés : analyse de variance
  séparée, sans prévision de direction. Aucun modèle n'est réestimé pour cette figure.}
\end{figure}

\FloatBarrier

\subsection{Règles de décision et comparaison R/Python}
Les règles autorisent SARIMA, VECM ou un ordre GARCH supérieur uniquement si
leurs gates scientifiques sont ouverts. L'absence d'estimation quand la
condition n'est pas soutenue est une décision documentée, pas un résultat
manquant à inventer. Les conventions et tables critiques natives R/Python
restent explicites, même lorsqu'elles diffèrent.

Le remplacement final de Naive exige simultanément RMSE et MAE non pires aux
deux horizons, au moins 1\,\% de gain sur la moyenne des ratios RMSE et le soutien
DM/Holm pré-déclaré, sans horizon significativement pire. Les résultats TEST
ne permettent aucun retuning. Non-rejet et équivalence, corrélation et causalité,
variance et direction du prix restent distingués.

\sourcecaption{Sources internes : \sourcepath{docs/protocol.md},
\sourcepath{docs/decisions.md}, protocoles gelés des Phases 3B, 4, 5, 6 et 7.
Le schéma est une synthèse documentaire ; il ne crée aucun résultat scientifique.}
